% Meeting Notes and TODOs — MT GitHub Mining Study
% Meeting date: 07/03/2026
% Project 43: Understanding the Adoption of Metamorphic Testing in Software Development
% Supervisor feedback recorded here for action tracking and inclusion in the P4P report.

\documentclass[12pt,a4paper]{article}
\usepackage[margin=2.5cm]{geometry}
\usepackage{hyperref}
\usepackage{enumitem}
\usepackage{booktabs}
\usepackage{xcolor}
\usepackage{mdframed}

\definecolor{todocolor}{RGB}{180,0,0}
\definecolor{donecolor}{RGB}{0,120,0}
\definecolor{opencolor}{RGB}{0,60,160}

\newenvironment{todo}{\begin{mdframed}[linecolor=todocolor,linewidth=1pt]\textbf{\color{todocolor}TODO: }}{\end{mdframed}}
\newenvironment{open}{\begin{mdframed}[linecolor=opencolor,linewidth=1pt]\textbf{\color{opencolor}OPEN QUESTION: }}{\end{mdframed}}
\newenvironment{done}{\begin{mdframed}[linecolor=donecolor,linewidth=1pt]\textbf{\color{donecolor}DONE: }}{\end{mdframed}}

\title{Meeting Notes and Research TODOs\\
       \large MT GitHub Mining Study --- Project 43}
\author{Ky Khon Quach (kqua446)}
\date{Meeting: 7 March 2026}

\begin{document}

\maketitle
\tableofcontents
\newpage

% ---------------------------------------------------------------------------
\section{Feedback from Meeting (07/03/2026)}
% ---------------------------------------------------------------------------

\subsection{Throughput Calculation --- Assumption Correction}

\begin{todo}
The throughput estimates in the previous report assumed fewer than 500 results
per query. This assumption is \textbf{incorrect}. The GitHub search API returns
up to 1,000 results per query (hard cap), and many MT-related queries are
expected to substantially exceed 500 results, particularly after the keyword
expansion described in Section~\ref{sec:keywords}. The throughput calculation
and any rate-limit timing arguments that depend on a ``$<$500 results''
assumption must be revised before the next meeting.
\end{todo}

\noindent
\textbf{Implication:} The adaptive date-range slicing mechanism in
\texttt{search.py} (activated when \texttt{total\_count > 900}) becomes more
important than previously estimated, since a higher fraction of queries will
require it. Time estimates for a full run should be recalculated accordingly.

% ---------------------------------------------------------------------------
\subsection{Keyword Expansion}
\label{sec:keywords}
% ---------------------------------------------------------------------------

The original keyword set covered only the canonical singular forms of core MT
terminology. The following variants must be added across \emph{all} five search
surfaces (repositories, code, issues/PRs, commits, and discussions):

\begin{itemize}
  \item \textbf{Plural forms:}
    \texttt{"metamorphic relations"}, \texttt{"metamorphic tests"},
    \texttt{"metamorphic relationships"}
  \item \textbf{Singular forms (where missing):}
    \texttt{"metamorphic test"} (no ``-ing'')
  \item \textbf{Synonym / related:}
    \texttt{"metamorphic relationship"}, \texttt{metamorphic\_test},
    \texttt{metamorphic\_relation} (underscore forms for code search)
\end{itemize}

\begin{done}
All of the above variants have been added to \texttt{config.py} for
\texttt{REPO\_QUERIES}, \texttt{CODE\_QUERIES}, \texttt{ISSUE\_QUERIES},
\texttt{COMMIT\_QUERIES}, and \texttt{DISCUSSION\_QUERIES}.
\end{done}

% ---------------------------------------------------------------------------
\subsection{File Extension Filter}
% ---------------------------------------------------------------------------

The original \texttt{ANALYSIS\_EXTENSIONS} set in \texttt{config.py} restricted
local analysis (Phase~3) to a fixed list of source-code extensions
(\texttt{.py}, \texttt{.java}, \texttt{.js}, etc.). This misses documentation
files (\texttt{.md}, \texttt{.rst}, \texttt{.txt}), configuration files
(\texttt{.yaml}, \texttt{.toml}), and notebook files (\texttt{.ipynb}) that may
contain explicit MT references (e.g., a \texttt{README.md} describing an MR).

\begin{done}
\texttt{ANALYSIS\_EXTENSIONS} has been changed to an empty set, which is
interpreted as ``accept all file extensions'' in \texttt{analyze.py}. To
restrict to specific extensions, populate the set manually in
\texttt{config.py}.
\end{done}

% ---------------------------------------------------------------------------
\subsection{Show Matched Query for All Result Types}
% ---------------------------------------------------------------------------

Previously, only Discussion records included a \texttt{matched\_query} field
indicating which search query produced the result. Repository, code, issue/PR,
and commit records did not carry this information, making it impossible to
trace a result back to the search term that found it.

\begin{done}
A \texttt{matched\_query} and \texttt{matched\_label} field is now added to
every record in \texttt{search.py} immediately after each query's
\texttt{\_sliced\_search} or \texttt{\_fetch\_all\_pages} call, for all five
search surfaces.
\end{done}

% ---------------------------------------------------------------------------
\subsection{Implicit Search Scope}
% ---------------------------------------------------------------------------

The original analysis default (\texttt{--only both}) included both explicit
and implicit pattern matching. The \texttt{IMPLICIT\_PATTERNS} defined in
\texttt{config.py} are intentionally broad structural heuristics (e.g., two
\texttt{assert} calls in close proximity) that will match the overwhelming
majority of any test file, regardless of whether the developer had any
awareness of MT. This creates serious noise in the results and directly
conflicts with the study's stated focus on \emph{explicit} MT adoption.

Specifically, the \texttt{follow\_up\_test\_call} pattern
(\texttt{(assert|assertEqual|expect).*\textbackslash n.*\textbackslash n.*(assert|assertEqual|expect)})
will fire on nearly every test file, producing thousands of false positives and
making the implicit hit counts meaningless as a signal of MT-like reasoning.

\begin{done}
The default value of \texttt{--only} in \texttt{analyze.py} has been changed
from \texttt{both} to \texttt{explicit}. Implicit analysis remains available
via \texttt{--only implicit} or \texttt{--only both}, but is no longer the
default. This aligns the tool with the study's focus on deliberate, explicit MT
adoption.
\end{done}

\noindent
\textbf{Longer-term:} The implicit patterns should either be substantially
tightened (e.g., require multiple assertions \emph{with an intervening input
transformation call}) or removed from the tool entirely if the study does not
plan to analyse implicit MT-like behaviour.

% ---------------------------------------------------------------------------
\subsection{General Code Clean-Up}
% ---------------------------------------------------------------------------

A general code review and clean-up pass is needed before the next meeting.
Specific items:

\begin{itemize}
  \item Remove or consolidate any duplicated logic between \texttt{search.py}
        and \texttt{config.py}.
  \item Ensure all output JSON files include a consistent schema (all records
        have \texttt{matched\_query} and \texttt{matched\_label}).
  \item Verify that deduplication logic in \texttt{search.py} handles the new
        keyword variants correctly (records from overlapping queries should
        not be double-counted).
  \item Add inline comments where non-obvious decisions are made (especially
        in \texttt{\_sliced\_search} and the scoring logic).
\end{itemize}

% ---------------------------------------------------------------------------
\subsection{Open Question --- Separation of Keywords Between Surfaces}
% ---------------------------------------------------------------------------

\begin{open}
Ask Ming Ray: why is there a separation of keywords between the different
search surfaces? For example, some terms appear in \texttt{REPO\_QUERIES} but
not in \texttt{CODE\_QUERIES}, and the \texttt{ISSUE\_QUERIES} set differs from
\texttt{COMMIT\_QUERIES}. Was this an intentional design decision (e.g., to
avoid redundancy where one surface subsumes another) or an oversight? The
answer affects whether we should unify the keyword sets or keep them surface-specific.
\end{open}

% ---------------------------------------------------------------------------
\section{IMPORTANT TODO: Step-by-Step Pipeline Description}
\label{sec:pipeline}
% ---------------------------------------------------------------------------

A clear, audience-facing description of the full mining pipeline is needed both
for the seminar presentation and for the P4P report methodology section. The
description should walk through each phase and define precisely what qualifies
as each tier of evidence.

\subsection{Phase 1 --- Search (search.py)}

\begin{enumerate}
  \item \textbf{Inputs:} GitHub API token, keyword lists in \texttt{config.py}.
  \item \textbf{Method:} For each keyword, query five search surfaces
        (repositories, code files, issues, pull requests, commit messages,
        GitHub Discussions) via the GitHub REST and GraphQL APIs.
  \item \textbf{Adaptive slicing:} If a query returns more than 900 results
        (GitHub hard cap is 1,000), the date window is bisected recursively
        (up to depth 6, i.e., 64 sub-windows) until each slice fits under the
        cap. This provides near-complete coverage.
  \item \textbf{Outputs:} Six deduplicated JSON files in \texttt{results/raw/}:
        \texttt{repos.json}, \texttt{code.json}, \texttt{issues.json},
        \texttt{prs.json}, \texttt{commits.json}, \texttt{discussions.json},
        plus a \texttt{summary.json} with per-query hit counts.
  \item \textbf{Key decision point:} A result at this stage means the keyword
        appeared \emph{somewhere} in the indexed surface (repo name/description,
        file content, issue body, commit message). It does not yet indicate
        deliberate MT adoption.
\end{enumerate}

\subsection{Phase 2 --- Clone (clone.py)}

\begin{enumerate}
  \item \textbf{Inputs:} \texttt{repos.json} and/or \texttt{code.json} from
        Phase~1.
  \item \textbf{Filtering:} Candidates are filtered by minimum stars, fork
        status, primary language, and maximum repository size (200~MB cap).
        Sorted by star count (descending).
  \item \textbf{Cloning:} Each candidate is shallow-cloned (depth=1, latest
        commit only) into \texttt{repos/<owner>/<name>/} using GitPython.
  \item \textbf{Outputs:} Cloned repository directories; a
        \texttt{results/cloned/manifest.json} tracking status of each repo
        (\texttt{cloned}, \texttt{skipped}, \texttt{refreshed}, \texttt{error}).
\end{enumerate}

\subsection{Phase 3 --- Analyse (analyze.py)}

\begin{enumerate}
  \item \textbf{Inputs:} Cloned repositories listed in the clone manifest.
  \item \textbf{Method:} Walk every file in each repo (all extensions by
        default). For each file, run explicit patterns. Optionally also run
        implicit patterns on test files only (\texttt{--only implicit} or
        \texttt{--only both}).
  \item \textbf{Outputs:}
    \begin{itemize}
      \item \texttt{findings.json} --- per-file findings with line number,
            snippet, pattern ID, and kind (explicit/implicit).
      \item \texttt{repo\_scores.json} --- per-repo aggregated scores and
            classification.
      \item \texttt{summary\_report.md} --- human-readable Markdown summary.
    \end{itemize}
\end{enumerate}

\subsection{Classification Tiers}
\label{sec:classification}

The following classification tiers define the strength of evidence for deliberate
MT adoption. These should be defined clearly in the P4P report and presented to
the seminar audience:

\begin{description}
  \item[\texttt{strong\_explicit}] ($\geq$3 explicit pattern hits in the repo).
    The repository contains multiple direct, unambiguous signals of MT adoption:
    e.g., imports of an MT library, use of a class or decorator named after MT,
    or multiple string/comment references to ``metamorphic testing'' or
    ``metamorphic relation''. This is the strongest evidence of deliberate
    adoption.

  \item[\texttt{weak\_explicit}] (1--2 explicit hits). At least one direct signal
    is present but it is isolated. The developer was likely aware of MT, but
    adoption may be minimal, exploratory, or limited to a single file or test.

  \item[\texttt{strong\_implicit}] ($\geq$5 implicit heuristic hits, no explicit
    hits). Structural patterns consistent with MT-like reasoning appear
    frequently in test files, but no explicit MT terminology is used. The
    developer may be applying MT intuitively without awareness of the technique.
    \textbf{Note:} The current implicit patterns are too broad and may produce
    false positives; these results require manual review.

  \item[\texttt{weak\_implicit}] (1--4 implicit hits, no explicit). A small number
    of MT-like structural patterns appear in test files. Low confidence; likely
    coincidental.

  \item[\texttt{none}] No signal found. The repository appeared in search results
    due to the keyword appearing in a README, commit message, or issue, but
    the source code itself contains no MT-related patterns.
\end{description}

% ---------------------------------------------------------------------------
\section{IMPORTANT TODO: Research Questions}
\label{sec:rqs}
% ---------------------------------------------------------------------------

The main research question for the mining component is:

\begin{quote}
  \textbf{RQ2 (mining): How prevalent is explicit metamorphic testing adoption
  across open-source repositories, and what project and author characteristics
  are associated with its presence?}
\end{quote}

The following sub-questions can be answered directly from the mined data:

\begin{enumerate}[label=\textbf{RQ2.\arabic*}]
  \item \textbf{Prevalence:} What proportion of repositories that mention MT
        (in any surface) actually contain explicit MT adoption in source code?
        (Answers the gap between ``awareness'' and ``adoption''.)

  \item \textbf{Trend over time:} Is the number of MT-adopting repositories
        growing, plateauing, or declining year-on-year? (Use
        \texttt{created\_at} from \texttt{repos.json}.)

  \item \textbf{Language distribution:} Which programming languages are most
        associated with explicit MT adoption? Does Python or Java dominate, as
        suggested by the library ecosystem?

  \item \textbf{Domain / repo type:} Are MT-adopting repositories concentrated
        in scientific computing, machine learning, compilers, or is adoption
        spread across domains? (Use repo topics, descriptions, and language as
        proxies for domain.)

  \item \textbf{Project maturity:} Are MT-adopting repositories more likely to
        be active production projects (high stars, active commit history) or
        research/prototype repositories? (Use stars, fork count, last-push date.)

  \item \textbf{Author concentration:} Are there specific authors, research
        groups, or organisations responsible for a disproportionate share of
        MT adoption? (Use author fields from commit records and issue records.)

  \item \textbf{Coverage within a repo:} How deeply does MT penetrate an
        adopting repository? Is it confined to a single test file or spread
        across multiple modules? (Use \texttt{affected\_files} from
        \texttt{repo\_scores.json}.)

  \item \textbf{MR structure:} What patterns do the explicit MR implementations
        follow? (Requires manual qualitative inspection of high-scoring
        repositories from \texttt{strong\_explicit} tier.)
\end{enumerate}

% ---------------------------------------------------------------------------
\section{IMPORTANT TODO: Results Report Structure}
\label{sec:results}
% ---------------------------------------------------------------------------

The results report should be organised to directly answer the sub-questions
above. Suggested structure:

\begin{enumerate}
  \item \textbf{Overview / funnel statistics} --- How many repos were found in
        search, how many were cloned, how many have explicit MT signal? Present
        as a funnel diagram.

  \item \textbf{Prevalence and classification breakdown} --- Table and/or bar
        chart of repos by classification tier
        (strong explicit / weak explicit / strong implicit / weak implicit / none).

  \item \textbf{Yearly trend of MT adoption} --- Line chart of number of
        MT-adopting repos created per year (use \texttt{created\_at}).
        Identifies whether MT uptake is accelerating or stagnant.

  \item \textbf{Language distribution} --- Bar chart of explicit MT repos by
        primary language. Highlights which language communities have adopted MT.

  \item \textbf{Domain analysis} --- Categorise high-scoring repos by domain
        (ML, scientific computing, web, systems, etc.) using topics and
        descriptions. Answers whether MT is confined to its original scientific
        computing niche.

  \item \textbf{Repository popularity / maturity} --- Scatter plot of
        classification tier vs.\ star count. Distinguishes research prototypes
        from production-quality projects.

  \item \textbf{Author / organisation concentration} --- Table of top authors
        and organisations by number of MT-adopting repos. Identifies MT
        ``champions'' and research-group clusters.

  \item \textbf{Pattern hit distribution} --- Which explicit patterns are most
        common? (From \texttt{summary\_report.md} pattern table.) Shows
        \emph{how} MT is expressed in code (imports, comments, function names).

  \item \textbf{Qualitative case studies} --- Manual inspection of the top
        5--10 \texttt{strong\_explicit} repos to describe the structure of
        real-world MR implementations (see Section~\ref{sec:qualitative}).
\end{enumerate}

\subsection{Suggested Plots}

\begin{itemize}
  \item Funnel chart: Search hits $\to$ repos cloned $\to$ explicit MT signal
  \item Bar chart: Classification tier distribution
  \item Line chart: MT-adopting repos created per year (2010--2025)
  \item Bar chart: Explicit MT repos by primary language (top 10 languages)
  \item Scatter plot: Stars vs.\ explicit hit count (log scale on x-axis)
  \item Heatmap: Language $\times$ Classification tier counts
  \item Bar chart: Pattern ID hit frequency (which patterns fire most)
\end{itemize}

% ---------------------------------------------------------------------------
\section{Qualitative Analysis}
\label{sec:qualitative}
% ---------------------------------------------------------------------------

In addition to the quantitative mining results, a qualitative dimension can
substantially enrich the study's contribution:

\subsection{Manual Inspection of High-Scoring Repos}

For the top 10--20 \texttt{strong\_explicit} repositories, conduct a manual
code review to answer:
\begin{itemize}
  \item What does the MR look like in practice? (How are the source test case,
        the follow-up test case, and the relational assertion structured?)
  \item Does the developer use a dedicated MT library or implement MRs manually?
  \item How many distinct MRs does the project define?
  \item What is the subject under test (e.g., a function, an API endpoint, a
        model prediction)?
\end{itemize}

\subsection{Discussion / Issue Mining}

The \texttt{discussions.json} and \texttt{issues.json} outputs can support a
lightweight qualitative analysis of \emph{how developers talk about MT}:
\begin{itemize}
  \item What questions do developers ask about MT in issues and discussions?
  \item What frustrations or barriers are mentioned?
  \item Are there requests for MT support, tool integration, or documentation?
\end{itemize}
This provides a complement to the quantitative prevalence data and directly
supports RQ1 (perceptual barriers) from the survey component.

\subsection{Sentiment Analysis (Optional)}

A lightweight sentiment analysis pass over the body text of issues and
discussions mentioning MT could reveal whether community discussions are
positive (``this technique is useful'', ``worked well for our project'') or
negative (``too hard to apply'', ``not worth the effort''). This would be a
novel empirical contribution and could be incorporated into the qualitative
findings.

% ---------------------------------------------------------------------------
\section{IMPORTANT TODO: Updated Results for Next Meeting}
% ---------------------------------------------------------------------------

\begin{todo}
Prepare an updated set of results for the next supervisor meeting. This should
include:
\begin{enumerate}
  \item A completed run of \texttt{search.py} with the expanded keyword set
        (Section~\ref{sec:keywords}), showing per-query hit counts from
        \texttt{summary.json}.
  \item A sample of cloned and analysed repos (even if not the full corpus),
        showing the classification distribution and top-scoring repos.
  \item At least one draft version of the yearly trend plot and the language
        distribution plot (Section~\ref{sec:results}).
  \item A revised throughput estimate that correctly accounts for the actual
        number of results returned per query (not the $<$500 assumption).
\end{enumerate}
\end{todo}

\end{document}
